\section{Análisis de catálogos y selección de componentes}
\label{sec:componentes}
% Fuentes iniciales sugeridas: Furukawa/Lightera, Huawei, VSOL, KOC, Wiretech (Córdoba).
% Se pueden usar otros fabricantes, incorporando las fichas técnicas y la fuente.

\subsection{Cables ópticos}
\label{sec:cables}
% Consigna 2.1: comparar cables para troncal/alimentación, distribución y acceso al abonado;
% tendido aéreo, por ductos o directamente enterrado.
% Para cada alternativa relevar: tipo de fibra y norma, cantidad de fibras y organización interna,
% atenuación por longitud de onda, construcción y protección, radio mínimo de curvatura,
% tracción máxima y demás parámetros mecánicos.
% Explicar qué cable se usa en cada tramo y justificar.
\pendiente{tabla comparativa de cables por tramo y justificación de la elección.}

\subsection{Componentes pasivos y equipos activos}
\label{sec:pasivos_activos}
% Consigna 2.2: cajas de empalme y cierres, cajas de distribución y terminales de acceso (NAP),
% splitters, conectores/adaptadores/pigtails/patch cords, distribuidores ópticos (ODF),
% OLT y ONU/ONT (los activos deben ser GPON).
% Presentar una tabla con: fabricante, modelo, función, especificaciones relevantes y aplicación prevista.
\pendiente{tabla de componentes pasivos y equipos activos GPON seleccionados.}
